% !TeX root = ../papers/nc-p3.tex
\section{二項四テンソルの整な種数一対応と整合する割線鉛筆}
\label{sec:elliptic-binary-cubic-pencils}
前節の有理正規三次 $C$ が点グラフの一次元の台全体で、残りの構造は有限とする。
各埋込点の埋込み次元は2以下とする（3の場合の判定は定理\ref{thm:maximal-tangent-ore}）。
四移動の曲線上の座標変換を $S$ と書き、四テンソルから四つの隣接角括弧を除いた多重一次式を $f$ とする。
\begin{theorem}\label{thm:elliptic-binary-cubic-pencils}
隣り合う二変数を固定し、残り二変数の $2\times2$ 行列式を取る。
このうち一つが整な $(2,2)$ 曲線なら、この窓は標準Koszul左右AS正則である。
節点型・尖点型も含み、各位置で高々四個の割線鉛筆の集合を選んで、隣り合う集合を関係で全射的に結ぶことができる。
\end{theorem}

\begin{proof}[順序付き二点の曲線]
四移動の継ぎ目を最後に置き、必要なら $S$ を逆に取る。
無限個の位置に各三次系を $C$ の座標で自明化する。
全ての隣接二座標で $f$ の評価行列が零になることはない。
対角外では元の四テンソルもその二点で消え、一意の点ファイバーに反する。
対角上の $(p,p)$ なら第二座標だけ曲線に沿って微分しても、括弧と $f$ が共に消えるので零となる。
これは第二位置の曲線の接線全体が第一位置の $p$ のファイバーに入ることを意味し、やはり不可能である。
$E$ を仮定の整因子とするとこの上で行列は階数1である。
左核、右核から、どの隣接二座標にも前と次の座標がそれぞれ一意に付く。
実際、隣接三点の条件は三変数を入れた四テンソルの商 $f$ が残る一変数で恒等零となることで、
隣の窓では同じ条件を巡回して読む。より正確には、任意の隣接行列式の零点スキームから、階数1行列の核を足して得る三点の
スキームは、隣の行列式の零点スキームからも逆側の核を足したものである。
これは単元成分で行列を掃き出せば階数1条件が残り一つの行列式となることによる。
従って隣の行列式が恒等零になることもなく、全て同じ整な算術種数1曲線となる。
一つの $E$ 上に全ての二重被覆 $p_j:E\to C$ を取れ、
$(p_{j-1},p_j)$ は対応する $(2,2)$ 曲線への同型である。
$D_j(x)=p_{j-1}(x)+p_j(x)\in\operatorname{Sym}^2C$ とする。
辺 $j$ には全ての $x\in E$ について次の割線関係がある。
\[
            D_j(x)\ \longmapsto\ S D_{j-2}(x).
\]
ここで因子 $d$ の鉛筆は $d S^1$ である。実際、四テンソルの先頭二変数に
$p_{j-2},p_{j-1}$ を入れると、残る二変数について $f$ は一次式の積で、
その根は $p_j,S p_{j-3}$ である。外側の二つの括弧が $p_{j-1},S p_{j-2}$ を加え、
中央の括弧が二鉛筆の階数2関係を与える。外側で評価した二点が一致する有限箇所にも
割線関係の閉性で延びる。四位置の移動で $D_{j+4}(E)=S D_j(E)$ となる。

\end{proof}
\begin{proof}[四次平面像の排除]
各辺の $\Pi=\mathbb P(T^\perp)$ は可逆な三行三列行列を含む。
実際、定理\ref{thm:twisted-cubic-residual}の基点の吹上げは完全交叉
$\{\det B=Q(B)=0\}$ に全射である。この像の計算は点グラフと無関係である。
$\Pi$ 上の吹上げの像で、例外因子以外から来るものは曲線外の有限個の点対から来る。
例外因子では $C$ の法微分を射影化しており、六式が消える法方向の次元は
点グラフの接空間の次元から1を引いたものである。従って有限個の埋込点の上で各高々一つの射影点しかない。
ゆえに $\Pi\cap\{\det B=Q=0\}$ は有限である。$\det B$ が $\Pi$ 上で恒等零なら
二次式だけで平面に無限個の零点を持つので矛盾する。

さらに $\Pi$ は片側の因子が二次Veronese曲線上にある階数1行列を含まない。
実際、基点写像 $B(X,Y)$ で $X$ に曲線上の $p$ を入れると、アフィン表示では
$B=u(p)(Y_{k+1}-pY_k)_k^t$、$u(p)=(1,p,p^2)^t$ を得る。
三差分は全射であるから、該当する非零行列が $\Pi$ にあれば $Y$ は $p$ 以外のファイバー点になる。
右側および無限遠点も同様である。

基底 $B_1,B_2,B_3$ を使い $M(d)$ の行を $d^t B_k$ とする。割線関係は $M(d)g=0$ である。
この三次行列式は恒等零にならず、階数1以下の軌跡は有限である。以下は $E$ の整性によらない。
可逆な $B_1$ を取れば、一般階数1以下は各 $d^t B_k$ が全 $d$ で $d^t B_1$ に比例することになり不可能である。
一般階数2で行列式が恒等零なら、随伴行列は多項式環の一意分解により $h(d)u(d)v(d)^t$ と書ける。
ここで $u,v$ は成分の最大公約数1の斉次ベクトルである（有理関数上で階数1に分解し分母を払う）。
定数左核は基底の独立性に反し、定数右核は可逆な $B_1$ に反する。
従って $u,v$ は一次である。$u$ が可逆な線型写像なら全 $d$ で $d\otimes u(d)\in T$ となり、
六次元の $T$ は対称二乗の片側を変換したものに等しい。直交する $\Pi$ に可逆行列がなく矛盾する。
そうでなければ、原始性により $u$ は階数2で、基底を変えて $u=(-d_2,d_1,0)^t$ とできる。
$M u=0$ より各行の先頭二成分は $(\lambda_k d_1,\lambda_k d_2)$ となる。
$\lambda$ の核から来る行列は一列しか持たず、$\Pi$ は階数1行列の直線を含む。
動く側の因子の直線は必ずVeronese二次曲線に交わるので、これも不可能である。
階数1以下の曲線があれば三次行列式の重複因子なので直線でなければならない。
その直線で $d^t B_k$ は $d^t B_1$ に比例し、$B_1$ の可逆性により各比例係数は直線上で一定である。
再び二つの独立な一次結合の行列はその左側二平面で消え、禁止された階数1行列の直線を得る。

$D_j(E)$ は二次曲線か四次曲線である。順序を忘れる射が次数2以下で
$D_j^*\mathcal O(1)$ の次数が4だからである。非零の三次行列式に入るため前者となり、
$D_j:E\to D_j(E)$ は次数2である。一般点で $M$ の階数が2なので $S D_{j-2}$ は有理的に $D_j$ を経由する。
両者が次数2なので被覆の対合 $s_j,s_{j-2}$ は等しい。
$s_j$ は $p_{j-1},p_j$ を交換する。順序付き二点の写像が埋込みだからである。
\end{proof}

\begin{proof}[楕円曲線の加法と鉛筆の連結]
滑らかな部分 $G$ に原点を取って群とする。滑らかな場合は楕円曲線、
節点なら $\mathbb G_m$、尖点なら $\mathbb G_a$ である。全て加法記号で書き、
$G\simeq\operatorname{Pic}^0 E$ は $x\mapsto\mathcal O(x-0)$ で同一視する。
この同一視は次数1の可逆層にRiemann--Rochで一意の零点があり、長さ1のCartier点は滑らかであることからも従う。
一つの二重被覆を $p$ とし、一般のファイバーの二点の和が一定なので、
原点を平行移動してその対合を $x\mapsto-x$ とする。$G$ 上の2倍は全射である。
$s_j$ は反射 $x\mapsto a_j-x$ である。実際、$G$ の自己同型は平行移動と原点を保つ群自己同型の合成で、
位数2ならその線型部は $1$ か $-1$ である。位数2の平行移動なら
$p_{j-1}$ の被覆対合と可換で、$p_j=p_{j-1}s_j$ は同じ対合を持ち、二点の埋込みに反する。
よって全 $j$ で $p_j(x)=p(x+c_j)$ と書ける。
交換対合の一致から $c_{j-1}+c_j=c_{j-3}+c_{j-2}$ が成り立つ。
また $L_j=p_j^*\mathcal O_C(1)$ とすると $L_jL_{j+3}=L_{j+1}L_{j+2}$ である。
実際、中二点で評価した階数1行列は $\mathbb P^1\times\mathbb P^1$ のSegre曲面への射を与え、
基点のないその四切断の束は $L_{j+1}L_{j+2}$、両射影で得る核の束は $L_j,L_{j+3}$ である。
$\operatorname{Pic}^2 E$ の加法座標で $L_j$ は $-2c_j$ なので、
$\delta_j=c_j-c_{j-1}$ とすれば $2\delta_{j+2}=2\delta_j$ となる。
さらに二周期の和 $c_{j-1}+c_j$ から、$c_{j+2}-c_j=t_j$ は $t_{j+1}=-t_j$ を満たし
$\delta_{j+2}-\delta_j=2t_j$ である。従って $4t_j=0$、$\delta_{j+4}=\delta_j$ である。

$S$ は $G$ に $v\mapsto \alpha v+b$ と持ち上がる。実際、四移動で隣接順序対の曲線を同型に写せる。
第一射影はそれぞれ $p$ を平行移動したものなので、これら平行移動で補正すれば $p$ に関する持上げとなる。
被覆対合との可換性から $2b=0$ である。
ここで $\alpha$ は原点を保つ自己同型であり、特殊な $j$ 不変量も許す。
二次曲線 $D_j(E)$ は順序を付けると $(p(v),p(v+\delta_j))$ の像で、これが決める差は符号を除き一意である。
実際、第一座標の一般持上げは $v,-v$ のみで、第二座標の一致で差を決定できる。
四移動の像と $\delta_{j+4}=\delta_j$ により $\alpha\delta_j=\pm\delta_j$ である。

$2d_j=\delta_j$ とし $m_j=c_{j-1}+d_j$ と取る。
$P(y,d)=p(y-d)+p(y+d)$ と書くと $D_j(x)=P(x+m_j,d_j)$ であり
\[
 P(y,d)=P(d,y),\qquad P(y+e,d)=P(y,d+e)\quad(e\in G[2]).
\]
各引数は独立に符号を変えられる。$2(m_j-m_{j-2})=0$ なので辺の対応は
$P(y,d_j)$ を $P(\alpha y+e_j,\alpha d_{j-2})$ に運ぶ（$e_j\in G[2]$）。
$\alpha d_j$ と $\pm d_j$ の差、$\alpha d_{j-2}$ と $\pm d_{j+2}$ の差はいずれも二捩れである。
従って有限集合
\[
 \{\,P(d_{j+1}+e,d_j):e\in G[2]\,\}
\]
から次の同じ集合に割線関係で全射する。実際、二引数を交換して二捩れを第一引数に集めれば
$P(d_{j+2}+e',d_{j+1})$ となり、$e$ を変えれば $e'$ が二捩れ全体を走る。
この列を両方向へ延ばし、系\ref{cor:twisted-cubic-compatible-pencils}を適用する。
\end{proof}

\begin{corollary}\label{cor:binary-vertical-ternary}
同じ窓で隣接行列式が恒等零、または一つでも垂直・水平な成分を持つ場合も標準正則である。
従ってこの三次曲線の残余でここまでの試験の外に残る二項対応は、
二本の $(1,1)$ グラフの和またはその二重化だけである。
\end{corollary}
\begin{proof}
一変数をその成分の値に固定する。次の変数を動かすと残り二位置の行列は全て階数1である。
この線型鉛筆には零行列もなく、Segre曲面の直線だから、どちらかの因子で核が一定である。
隣接位置の一定核なら先に排除した隣接零行列になるため、反対位置で一定核となる。
従って四テンソルも相対する一対の曲線点で消える。
一方の点での三テンソルの中支持は高々3である。両端支持は点ファイバー条件で3、
各 $R_i$ に階数1関係はない。定理\ref{thm:finite-frobenius-low-opposite-as}の末段の鉛筆の議論により
中支持は3であり、定理\ref{thm:ternary-point-subalgebras}の三方向の縮約にも階数1はない。
これで全次数正則性を得る。残る因子の主張は二次数を加えるだけである。
\end{proof}
